% ============================================================
%  РАЗДЕЛ 2.3  Разработка макетов приложения
% ============================================================

\section{Разработка макетов приложения}

Проектирование пользовательского интерфейса является самостоятельным этапом разработки программного комплекса. Высокодетализированные прототипы экранных форм позволяют верифицировать логику навигации, выявить противоречия в расположении элементов управления и зафиксировать визуальное соглашение до начала реализации HTML-шаблонов. В соответствии с методическими требованиями в данном разделе обосновывается выбор цветового решения, типографики и стиля интерфейса, а также приводятся прототипы всех ключевых экранных форм разрабатываемой системы.

\subsection{Обоснование стиля и цветового решения}

Целевая аудитория системы --- инженеры тестирования, работающие в условиях серверных помещений и рабочих станций с профессиональными мониторами. Эксплуатация системы носит интенсивный характер: оператор одновременно отслеживает состояние нескольких десятков заглушек, переключает тумблеры Rate Limiter, запускает и останавливает процессы. Это определяет два ключевых требования к интерфейсу: минимальную утомляемость при длительной работе и мгновенную читаемость статусов без необходимости анализировать текстовые подписи.

Для выполнения первого требования выбрана тёмная цветовая схема (dark theme): светлый текст на тёмном фоне снижает нагрузку на зрение при многочасовой работе и является де-факто стандартом для инструментов класса DevOps и Operations. Для выполнения второго требования разработана система цветовых статусных индикаторов, в которой цвет однозначно кодирует состояние объекта: зелёный соответствует активному состоянию (\texttt{RUNNING}), жёлтый~--- промежуточному (\texttt{REGISTERED}~/~\texttt{STOPPED}), красный~--- ошибочному (\texttt{ERROR}~/~\texttt{UNAVAILABLE}).

Акцентный цвет интерфейса~--- синий (\texttt{\#4F8EF7}). Данный оттенок обеспечивает достаточный контраст по отношению к тёмному фону (коэффициент контрастности превышает 4,5:1, что соответствует требованиям WCAG~2.1 уровня~AA) и не несёт эмоциональной нагрузки, характерной для статусных цветов. Акцентным цветом выделяются интерактивные элементы: кнопки основного действия, активные пункты навигации и ссылки.

Типографическую основу интерфейса составляет шрифт Inter, который входит в базовую поставку большинства дистрибутивов Linux и не требует загрузки с внешних CDN. Это соответствует требованию технологической независимости от зарубежных сервисов, сформулированному в разделе~1.3. Основной размер шрифта~--- 14~px, межстрочный интервал~--- 1.5. Для отображения технических идентификаторов (имена файлов, JVM-аргументы, пути к директориям) используется моноширинный шрифт JetBrains~Mono, обеспечивающий лучшую читаемость длинных строк технического содержания.

Компоновка страниц построена на двухколоночной сетке: постоянная боковая панель навигации шириной 220~px и основная рабочая область, занимающая оставшееся пространство. Такая компоновка является стандартной для административных веб-панелей и обеспечивает предсказуемость навигации при переключении между разделами системы. Сводная характеристика визуальной системы интерфейса приведена в таблице~\ref{tab:ui-colors}.

\begin{table}[H]
\caption{Визуальная система интерфейса}
\label{tab:ui-colors}
\centering
\begin{tblr}{
  colspec = {|l|l|X|},
  hlines,
  width = \linewidth,
}
Элемент & Значение & Назначение \\
Фон основного контента & \texttt{\#121218} & Снижение утомляемости при длительной работе \\
Фон боковой панели & \texttt{\#1A1A24} & Визуальное разделение навигации и контента \\
Акцентный цвет & \texttt{\#4F8EF7} & Интерактивные элементы, активные состояния \\
Статус RUNNING & \texttt{\#22C55E} (зелёный) & Процесс запущен и активен \\
Статус REGISTERED / STOPPED & \texttt{\#EAB308} (жёлтый) & Процесс зарегистрирован, не запущен \\
Статус ERROR / UNAVAILABLE & \texttt{\#EF4444} (красный) & Ошибка или недоступность хоста \\
Основной текст & \texttt{\#E2E8F0} & Высокий контраст на тёмном фоне \\
Вторичный текст & \texttt{\#94A3B8} & Метаданные, метки полей, подсказки \\
Основной шрифт & Inter & Читаемость, доступность в Linux \\
Моноширинный шрифт & JetBrains Mono & Технические идентификаторы и аргументы \\
\end{tblr}
\end{table}

\subsection{Структура навигации и общий хром приложения}

Все экранные формы системы разделяют единый шаблон-«обёртку», включающий постоянную боковую панель навигации и верхнюю строку с заголовком текущего раздела. Структура навигации включает три основных раздела:

\begin{itemize}
    \item \textbf{Dashboard}~--- главная панель управления заглушками; отображается по умолчанию при входе в систему;
    \item \textbf{Настройки}~--- управление хостами и учётными записями SSH, глобальные параметры системы;
    \item \textbf{Журналы}~--- просмотр вывода Java-процессов заглушек.
\end{itemize}

Боковая панель содержит логотип системы «MockControl» в верхней части и строку версии в нижней. Правый верхний угол заголовочной полосы показывает оперативный статус подключения к Redis, что позволяет оператору немедленно обнаруживать сбои хранилища, не выполняя диагностических действий. Макет общего хрома приложения с боковой навигацией представлен на рисунке~\ref{fig:mockup-shell}.

\begin{figure}[H]
\centering
\includegraphics[width=\textwidth]{mockup-shell.png}
\caption{Макет общего хрома приложения с боковой навигацией}
\label{fig:mockup-shell}
\end{figure}

%% ============================================================
%% ПРОМТ ДЛЯ FIGMA AI — рисунок «Общий хром приложения»
%% ============================================================
%%
%% Design a dark admin web application shell (combined wireframe and
%% visual design). Canvas: 1440×900 px. All UI labels in Russian.
%%
%% COLOR SYSTEM
%%   Page background ............. #121218
%%   Sidebar background ........... #1A1A24
%%   Accent / primary action ....... #4F8EF7
%%   Primary text .................. #E2E8F0
%%   Secondary text ................ #94A3B8
%%   Border / divider .............. #2D3748
%%   Card / panel background ....... #1E1E2E
%%
%% TYPOGRAPHY
%%   All text: Inter. Monospace accents: JetBrains Mono.
%%
%% LEFT SIDEBAR — fixed, 220 px wide, full height, background #1A1A24:
%%   • Top block (padding 24 px): hexagon icon "⬡" (20 px, #4F8EF7) +
%%     text "MockControl" (16 px bold, #E2E8F0) on the same line.
%%   • Navigation list starting 32 px below the logo, each item 44 px tall,
%%     full width, horizontal padding 16 px:
%%       Item 1 — icon LayoutDashboard + label "Dashboard"
%%         ACTIVE STATE: background rgba(79,142,247,0.12),
%%         left border 3 px solid #4F8EF7, icon and text color #4F8EF7.
%%       Item 2 — icon Settings + label "Настройки"
%%         Inactive: icon and text #94A3B8.
%%       Item 3 — icon ScrollText + label "Журналы"
%%         Inactive.
%%   • Bottom of sidebar (pinned): version tag "v 1.0.0"
%%     (12 px, #94A3B8, padding-bottom 16 px, centered).
%%
%% TOP HEADER BAR — spans full width minus sidebar, height 56 px,
%%   background #1A1A24, bottom border 1 px solid #2D3748:
%%   • Left (padding-left 24 px): text "Dashboard" (18 px semibold, #E2E8F0).
%%   • Right (padding-right 24 px): filled circle 8 px #22C55E + text
%%     "Redis: подключён" (13 px, #94A3B8), items vertically centered.
%%
%% MAIN CONTENT AREA — remainder, background #121218, padding 24 px:
%%   Show only an empty placeholder: a dashed rectangle with rounded corners
%%   (border 2 px dashed #2D3748, border-radius 12 px, full available width,
%%   height ~560 px) containing centered text "Область контента"
%%   (14 px, #2D3748).
%%
%% Style rules: no background gradients, 1 px borders, 8 px rounded corners
%%   on cards, subtle box-shadow 0 1 px 3 px rgba(0,0,0,0.4) on cards only.
%% ============================================================


\subsection{Макет главной панели управления (Dashboard)}

Dashboard является основным рабочим экраном системы. Он отображает полный список зарегистрированных заглушек и предоставляет оператору возможность управлять их жизненным циклом непосредственно из таблицы без перехода на отдельные страницы.

Верхняя часть рабочей области занята четырьмя информационными карточками со сводной статистикой: общее число зарегистрированных заглушек, число активных (\texttt{RUNNING}), число заглушек в статусе ошибки и число доступных хостов. Карточки позволяют оценить состояние всей инфраструктуры одним взглядом без анализа содержимого таблицы.

Таблица заглушек содержит столбцы: имя файла (уникальный идентификатор, отображается моноширинным шрифтом), целевой хост, статус процесса в виде цветного бейджа, назначенный порт, тумблер Rate Limiter и кнопки действий (Запустить~/ Остановить~/ Удалить). Кнопка «Зарегистрировать заглушку» в правом верхнем углу открывает модальное окно загрузки артефакта.

Макет Dashboard с заполненной таблицей представлен на рисунке~\ref{fig:mockup-dashboard}.

\begin{figure}[H]
\centering
\includegraphics[width=\textwidth]{mockup-dashboard.png}
\caption{Макет главной панели управления (Dashboard)}
\label{fig:mockup-dashboard}
\end{figure}

%% ============================================================
%% ПРОМТ ДЛЯ FIGMA AI — рисунок «Dashboard»
%% ============================================================
%%
%% Full Dashboard screen for a dark admin web app.
%% Canvas: 1440×900 px. Same sidebar + header shell as the previous
%% prompt. Active nav item: "Dashboard". All labels in Russian.
%% Color system and typography: identical to shell prompt.
%%
%% MAIN CONTENT AREA (padding 24 px):
%%
%% ── ROW 1: PAGE HEADING + ACTION BUTTON (height 40 px, flex, align center)
%%   Left:  "Заглушки" — 20 px semibold, #E2E8F0.
%%   Right: primary button "＋ Зарегистрировать заглушку"
%%          background #4F8EF7, text white 14 px medium, height 40 px,
%%          border-radius 8 px, horizontal padding 16 px.
%%
%% ── ROW 2: STATS CARDS (margin-top 20 px, 4 equal cards, gap 16 px)
%%   Each card: background #1E1E2E, border 1 px #2D3748, border-radius 12 px,
%%   padding 20 px, height ~96 px, flex column.
%%   Card structure: small icon top-left (#94A3B8, 18 px Lucide) +
%%                   label text below (13 px #94A3B8) +
%%                   large number at bottom (32 px bold).
%%
%%   Card 1 — icon Layers,       label "Всего заглушек",  number "7"  color #E2E8F0
%%   Card 2 — icon PlayCircle,   label "Запущено",        number "4"  color #22C55E
%%   Card 3 — icon AlertCircle,  label "Ошибок",          number "1"  color #EF4444
%%   Card 4 — icon Server,       label "Хостов доступно", number "3"  color #4F8EF7
%%
%% ── ROW 3: MOCKS TABLE (margin-top 24 px)
%%   Container: background #1E1E2E, border-radius 12 px, border 1 px #2D3748,
%%   overflow hidden.
%%
%%   TABLE HEADER ROW: background #161622, height 40 px, padding 0 20 px,
%%   text #94A3B8 12 px uppercase letter-spacing 0.05 em.
%%   Columns (px widths): "Имя файла" 260 | "Хост" 160 | "Статус" 120 |
%%   "Порт" 80 | "Rate Limiter" 130 | "Действия" 170
%%
%%   DATA ROWS — 5 rows, height 52 px each, horizontal padding 20 px,
%%   alternating background #1E1E2E / #1B1B2B, font 14 px.
%%
%%   Row 1: "payment-service-stub.jar" (JetBrains Mono 13 px, #E2E8F0) |
%%          "test-server-01" (#94A3B8 13 px) |
%%          badge "RUNNING" (pill, bg rgba(34,197,94,0.12), text #22C55E,
%%          border 1 px rgba(34,197,94,0.3), 12 px, padding 2 px 10 px) |
%%          "8101" (#E2E8F0) |
%%          toggle ON (#4F8EF7 filled pill 36×20 px, white circle right) |
%%          [Stop btn] [Delete btn]
%%
%%   Row 2: "auth-service-stub.jar" | "test-server-01" |
%%          badge "RUNNING" | "8102" | toggle OFF (grey) |
%%          [Stop] [Delete]
%%
%%   Row 3: "notification-stub.jar" | "test-server-02" |
%%          badge "STOPPED" (pill, bg rgba(234,179,8,0.12), text #EAB308,
%%          border 1 px rgba(234,179,8,0.3)) |
%%          "—" (#94A3B8) | toggle OFF disabled (opacity 0.4) |
%%          [Start btn] [Delete]
%%
%%   Row 4: "legacy-api-mock.jar" | "localhost" |
%%          badge "ERROR" (pill, bg rgba(239,68,68,0.12), text #EF4444,
%%          border 1 px rgba(239,68,68,0.3)) |
%%          "—" | toggle OFF disabled |
%%          [Start] [Delete]
%%
%%   Row 5: "billing-stub.jar" | "test-server-02" |
%%          badge "REGISTERED" (#EAB308 style) | "—" | toggle OFF | [Start] [Delete]
%%
%%   BUTTON STYLES (height 32 px, border-radius 6 px, 13 px, padding 0 12 px):
%%   "Стоп"   — border 1 px #94A3B8, text #94A3B8, bg transparent.
%%   "Старт"  — border 1 px #22C55E, text #22C55E, bg transparent.
%%   "Удалить"— icon Trash2 only, 16 px, color #94A3B8, bg transparent,
%%              hover color #EF4444.
%%
%%   TOGGLE: pill 36×20 px.
%%   ON: background #4F8EF7, white circle 16 px right-aligned.
%%   OFF: background #2D3748, circle #94A3B8 left-aligned.
%%   Disabled OFF: same as OFF at opacity 0.4.
%%
%% No modal overlay in this frame.
%% ============================================================


\subsection{Макет модального окна регистрации заглушки}

Регистрация новой заглушки выполняется через модальное окно, открывающееся поверх Dashboard без перехода на отдельную страницу. Такой подход сохраняет контекст: оператор видит уже зарегистрированные заглушки за полупрозрачным затемнённым фоном и может убедиться в отсутствии конфликта имён до подтверждения операции.

Форма содержит пять полей: зону перетаскивания файла артефакта (\texttt{.jar}~/~\texttt{.war}), выпадающий список целевого хоста, текстовое поле JVM-аргументов, числовое поле лимита запросов Rate Limiter и флажок немедленного запуска после регистрации. Поле JVM-аргументов отображается моноширинным шрифтом, что снижает вероятность ошибок при вводе длинных строк аргументов.

Макет Dashboard с открытым модальным окном регистрации представлен на рисунке~\ref{fig:mockup-register}.

\begin{figure}[H]
\centering
\includegraphics[width=\textwidth]{mockup-register.png}
\caption{Макет модального окна регистрации новой заглушки}
\label{fig:mockup-register}
\end{figure}

%% ============================================================
%% ПРОМТ ДЛЯ FIGMA AI — рисунок «Модальное окно регистрации»
%% ============================================================
%%
%% Same dashboard frame as the previous prompt, but with a modal overlay.
%% Canvas: 1440×900 px. Language: Russian.
%%
%% DIMMED BACKGROUND: the full dashboard (table + stats cards) is visible
%% but covered by a semi-transparent layer rgba(0,0,0,0.65).
%%
%% MODAL WINDOW — horizontally and vertically centered on canvas:
%%   Width: 520 px. Background: #1E1E2E. Border: 1 px solid #2D3748.
%%   Border-radius: 16 px. Padding: 32 px.
%%   Box-shadow: 0 24 px 64 px rgba(0,0,0,0.7).
%%
%%   HEADER ROW (flex, space-between):
%%     Left: "Регистрация заглушки" — 18 px semibold, #E2E8F0.
%%     Right: icon button X (20 px, color #94A3B8).
%%
%%   FORM — vertical stack, gap 20 px, margin-top 24 px:
%%
%%   Field 1 — FILE UPLOAD ZONE:
%%     Label "Артефакт (.jar / .war)" — 13 px, #94A3B8, mb 8 px.
%%     Dashed rectangle: border 2 px dashed #2D3748, border-radius 10 px,
%%     height 88 px, background #161622.
%%     Centered: icon Upload (24 px, #4F8EF7) above text
%%     "Перетащите файл или нажмите для выбора" (14 px, #94A3B8).
%%     Below that: ".jar, .war  •  до 512 МБ" (12 px, #64748B).
%%
%%   Field 2 — DROPDOWN "Целевой хост":
%%     Label 13 px #94A3B8.
%%     Select: full width, height 44 px, background #161622,
%%     border 1 px #2D3748, border-radius 8 px, padding 0 12 px.
%%     Placeholder "Выберите хост..." (#64748B, 14 px). Chevron-down right.
%%
%%   Field 3 — TEXT INPUT "JVM-аргументы":
%%     Label 13 px #94A3B8.
%%     Input: full width, height 44 px, same style.
%%     Placeholder "-Xms128m -Xmx256m" in JetBrains Mono 13 px, #64748B.
%%
%%   Field 4 — NUMBER INPUT "Лимит запросов (req/s)":
%%     Label 13 px #94A3B8.
%%     Input: width 160 px, height 44 px, same style. Value "500".
%%     Helper text below "0 — без ограничений" (12 px, #64748B).
%%
%%   Field 5 — CHECKBOX ROW:
%%     Checkbox (16×16 px, border 1 px #2D3748, border-radius 4 px,
%%     unchecked state, bg #161622) + label
%%     "Запустить сразу после регистрации" (14 px, #94A3B8).
%%
%%   FOOTER (margin-top 28 px, flex, justify-end, gap 12 px):
%%     "Отмена" — height 40 px, border 1 px #2D3748, bg transparent,
%%                text #94A3B8, border-radius 8 px, padding 0 20 px.
%%     "Зарегистрировать" — height 40 px, bg #4F8EF7, text white,
%%                border-radius 8 px, padding 0 20 px, 14 px medium.
%% ============================================================


\subsection{Макет страницы настроек: управление хостами}

Страница настроек организована в виде двух горизонтальных вкладок: <<Хосты>> и <<Учётные записи>>. Такое разделение предотвращает перегруженность экрана и группирует логически связанные объекты.

Вкладка <<Хосты>> отображает таблицу зарегистрированных целевых серверов с индикатором доступности SSH-соединения для каждого хоста. Столбец <<Проверено>> показывает дату и время последней проверки доступности, что позволяет оценить актуальность отображаемого статуса. Кнопка <<Добавить хост>> открывает модальную форму с полями, описанными в разделе~2.1.6.

Вкладка <<Учётные записи>> отображает список сохранённых учётных записей ОС Linux для SSH/SCP-операций. Пароли в таблице не отображаются: вместо них показывается замаскированная строка «\texttt{••••••••}», что исключает случайную утечку при демонстрации интерфейса коллегам или на презентации.

Макет страницы настроек (вкладка <<Хосты>>) представлен на рисунке~\ref{fig:mockup-settings}, форма добавления нового хоста~--- на рисунке~\ref{fig:mockup-add-host}.

\begin{figure}[H]
\centering
\includegraphics[width=\textwidth]{mockup-settings.png}
\caption{Макет страницы настроек, вкладка «Хосты»}
\label{fig:mockup-settings}
\end{figure}

%% ============================================================
%% ПРОМТ ДЛЯ FIGMA AI — рисунок «Настройки — вкладка Хосты»
%% ============================================================
%%
%% Settings page for a dark admin web app.
%% Canvas: 1440×900 px. Same shell. Active nav item: "Настройки".
%% Header title: "Настройки". All labels in Russian.
%% Color system: identical to all previous prompts.
%%
%% MAIN CONTENT (padding 24 px):
%%
%% ── ROW 1: PAGE HEADING
%%   "Настройки" — 20 px semibold, #E2E8F0.
%%
%% ── ROW 2: TAB BAR (margin-top 20 px)
%%   Two tabs, underline indicator style, gap 0, border-bottom 1 px #2D3748
%%   spanning full width:
%%     Tab "Хосты"         — ACTIVE: text #4F8EF7, 15 px medium,
%%                           border-bottom 2 px solid #4F8EF7, pb 12 px.
%%     Tab "Учётные записи"— inactive: text #94A3B8, same sizing.
%%   Tabs spaced 32 px apart.
%%
%% ── ROW 3: TABLE CONTAINER (margin-top 24 px)
%%   Background #1E1E2E, border-radius 12 px, border 1 px #2D3748,
%%   overflow hidden.
%%
%%   TOOLBAR (padding 16 px 20 px, flex, space-between, align center):
%%     Left: "Целевые хосты развёртывания" — 14 px, #94A3B8.
%%     Right: "＋ Добавить хост" — bg #4F8EF7, white text 14 px, height 40 px,
%%             border-radius 8 px, padding 0 16 px.
%%
%%   HEADER ROW: bg #161622, height 40 px, padding 0 20 px,
%%   text #94A3B8 12 px uppercase.
%%   Columns (px): "Хост / Адрес" 200 | "SSH-порт" 80 |
%%   "Рабочая директория" 210 | "Статус" 120 |
%%   "Проверено" 150 | "Описание" auto | "Действия" 90
%%
%%   DATA ROWS — 3 rows, height 52 px, padding 0 20 px, font 14 px,
%%   alternating bg #1E1E2E / #1B1B2B:
%%
%%   Row 1:
%%     "test-server-01" (JetBrains Mono 13 px, #E2E8F0) |
%%     "22" (#94A3B8) |
%%     "/opt/mock-services" (Mono 12 px, #94A3B8) |
%%     badge "AVAILABLE" (bg rgba(34,197,94,0.12), text #22C55E,
%%     border 1 px rgba(34,197,94,0.3), 12 px pill) |
%%     "01.11.2024 10:29" (13 px #94A3B8) |
%%     "Стенд нагрузочного тестирования" (13 px #94A3B8) |
%%     [pencil icon btn] [trash icon btn]
%%
%%   Row 2: "test-server-02" | "22" | "/opt/mock-services" |
%%     badge "AVAILABLE" | "01.11.2024 10:29" | "Контур препром" |
%%     [edit] [delete]
%%
%%   Row 3: "localhost" | "22" | "/opt/mock-services" |
%%     badge "UNAVAILABLE" (bg rgba(239,68,68,0.12), text #EF4444,
%%     border 1 px rgba(239,68,68,0.3)) |
%%     "01.11.2024 09:15" | "Локальный сервер управления" |
%%     [edit] [delete]
%%
%%   Icon buttons: pencil and trash2, 16 px, color #94A3B8,
%%   hover pencil → #4F8EF7, hover trash → #EF4444.
%%
%% No modal overlay in this frame.
%% ============================================================

\begin{figure}[H]
\centering
\includegraphics[width=\textwidth]{mockup-add-host.png}
\caption{Макет модального окна добавления нового хоста}
\label{fig:mockup-add-host}
\end{figure}

%% ============================================================
%% ПРОМТ ДЛЯ FIGMA AI — рисунок «Модальное окно — добавление хоста»
%% ============================================================
%%
%% Same Settings/Хосты page dimmed with rgba(0,0,0,0.65) overlay.
%% Canvas: 1440×900 px. Language: Russian.
%%
%% MODAL — centered:
%%   Width: 560 px. Background: #1E1E2E. Border: 1 px #2D3748.
%%   Border-radius: 16 px. Padding: 32 px.
%%   Box-shadow: 0 24 px 64 px rgba(0,0,0,0.7).
%%
%%   HEADER: "Добавить хост" (18 px semibold, #E2E8F0) + close X (#94A3B8).
%%
%%   FORM (vertical stack, gap 18 px, margin-top 24 px):
%%
%%   Field 1 "Hostname / IP-адрес":
%%     Full-width input, height 44 px, bg #161622, border 1 px #2D3748,
%%     border-radius 8 px. Value "test-server-03" (Mono 13 px, #E2E8F0).
%%
%%   Field 2 "SSH-порт":
%%     Width 120 px, same style. Value "22".
%%
%%   Field 3 "Учётная запись":
%%     Full-width dropdown. Selected "mock-runner" (#E2E8F0). Chevron right.
%%
%%   Field 4 "Рабочая директория":
%%     Full-width input. Value "/opt/mock-services" (Mono, #E2E8F0).
%%
%%   Field 5 "Путь к Java":
%%     Full-width input.
%%     Value "/usr/lib/jvm/java-17-openjdk-amd64/bin/java" (Mono 12 px, #E2E8F0).
%%     Helper: "Абсолютный путь к исполняемому файлу java" (12 px, #64748B).
%%
%%   Field 6 — TWO INPUTS IN ONE ROW (gap 16 px):
%%     Left: "Порт мин." label + input width ~220 px, value "8100".
%%     Right: "Порт макс." label + input width ~220 px, value "8200".
%%
%%   Field 7 "Описание (необязательно)":
%%     Textarea full width, height 72 px, bg #161622, border 1 px #2D3748,
%%     border-radius 8 px.
%%     Placeholder "Краткое описание назначения хоста..." (#64748B, 14 px).
%%
%%   FOOTER (mt 28 px, flex, justify-end, gap 12 px):
%%     "Отмена" — grey outline, height 40 px, border-radius 8 px.
%%     "Сохранить" — bg #4F8EF7, white, height 40 px, border-radius 8 px.
%% ============================================================


\subsection{Макет страницы просмотра журналов процессов}

Страница журналов предназначена для просмотра стандартного вывода (stdout/stderr) Java-процессов заглушек. Выбор заглушки осуществляется через выпадающий список, что позволяет переключаться между логами без перехода на отдельные страницы.

Основную часть экрана занимает терминальный блок с монохромным текстом на чёрном фоне. Строки, содержащие ключевые слова \texttt{ERROR} и
\texttt{WARN}, выделяются соответственно красным и жёлтым фоном для
быстрого визуального поиска проблем. Метки временных меток окрашены в нейтральный серо-синий цвет, чтобы не конкурировать с читаемостью основного текста.

Панель управления над терминальным блоком содержит: выпадающий список заглушек, тумблер автообновления, кнопку ручного обновления и кнопку очистки отображаемого буфера. Строка состояния между панелью управления и терминальным блоком показывает PID процесса и назначенный порт, что позволяет оператору проверить корректность отображаемых данных. Макет страницы журналов с заполненным выводом представлен на рисунке~\ref{fig:mockup-logs}.

\begin{figure}[H]
\centering
\includegraphics[width=\textwidth]{mockup-logs.png}
\caption{Макет страницы просмотра журналов процессов}
\label{fig:mockup-logs}
\end{figure}

%% ============================================================
%% ПРОМТ ДЛЯ FIGMA AI — рисунок «Журналы процессов»
%% ============================================================
%%
%% Logs page for a dark admin web app.
%% Canvas: 1440×900 px. Same shell. Active nav item: "Журналы".
%% Header title: "Журналы процессов". Language: Russian.
%% Color system: identical to all previous prompts.
%%
%% MAIN CONTENT (padding 24 px):
%%
%% ── ROW 1: CONTROLS BAR (flex, align center, gap 12 px, height 44 px)
%%
%%   Left group (flex, align center, gap 8 px):
%%     Label "Заглушка:" (14 px, #94A3B8)
%%     Dropdown (width 300 px, height 40 px, bg #1E1E2E,
%%     border 1 px #2D3748, border-radius 8 px, px 12 px):
%%       Selected value "payment-service-stub.jar"
%%       (JetBrains Mono 13 px, #E2E8F0). Chevron-down right.
%%
%%   Right group (margin-left auto, flex, align center, gap 12 px):
%%     Toggle pill 36×20 px (ON, #4F8EF7) + label "Авто" (13 px, #94A3B8)
%%     Button "↻ Обновить" (border 1 px #2D3748, bg transparent,
%%     text #94A3B8, 13 px, height 36 px, br 8 px, px 14 px)
%%     Button "Очистить" (border 1 px rgba(239,68,68,0.4), bg transparent,
%%     text #EF4444, same size)
%%
%% ── ROW 2: STATUS BAR (margin-top 16 px, height 36 px,
%%   bg #1A1A24, border-radius 8 px 8 px 0 0,
%%   border 1 px #2D3748 on top+sides, padding 0 16 px,
%%   flex, align center, space-between):
%%   Left: dot 6 px #22C55E + text
%%   "payment-service-stub.jar  —  PID 24731  —  порт 8101"
%%   (JetBrains Mono 12 px, #94A3B8).
%%   Right: "Строк в буфере: 847" (12 px, #64748B).
%%
%% ── ROW 3: TERMINAL BLOCK (flex-grow, min-height 540 px,
%%   bg #0D0D14, border-radius 0 0 8 px 8 px,
%%   border 1 px #2D3748 on bottom+sides,
%%   padding 16 px, overflow-y auto)
%%
%%   Font: JetBrains Mono 12.5 px, line-height 20 px.
%%   Show 16 log lines. Format per line:
%%     timestamp (color #4A5568) + level keyword + logger name + message.
%%
%%   NORMAL lines (text #C8D3E0):
%%   "2024-11-01 10:35:02.114  INFO  24731 --- [main] o.s.b.SpringApplication     : Starting PaymentStubApplication"
%%   "2024-11-01 10:35:02.891  INFO  24731 --- [main] o.s.b.SpringApplication     : No active profile set, falling back to 1 default profile"
%%   "2024-11-01 10:35:04.210  INFO  24731 --- [main] o.s.web.TomcatWebServer     : Tomcat initialized with port 8101"
%%   "2024-11-01 10:35:05.003  INFO  24731 --- [main] o.s.b.SpringApplication     : Started in 2.889 seconds"
%%   "2024-11-01 10:35:15.441  INFO  24731 --- [exec-1] c.e.stub.PaymentController  : GET /payment/status → 200 OK"
%%   "2024-11-01 10:35:20.112  INFO  24731 --- [exec-2] c.e.stub.PaymentController  : POST /payment/process → 200 OK"
%%
%%   WARN line (full-row bg rgba(234,179,8,0.08)):
%%   "2024-11-01 10:35:22.778  WARN  24731 --- [exec-3] c.e.stub.PaymentController  : Slow response detected: 843 ms"
%%   keyword "WARN" colored #EAB308, rest of line #EAB308.
%%
%%   NORMAL lines continue:
%%   "2024-11-01 10:35:30.001  INFO  24731 --- [exec-4] c.e.stub.PaymentController  : GET /payment/limits → 200 OK"
%%   "2024-11-01 10:35:38.220  INFO  24731 --- [exec-5] c.e.stub.PaymentController  : POST /payment/process → 200 OK"
%%
%%   ERROR line (full-row bg rgba(239,68,68,0.10)):
%%   "2024-11-01 10:35:45.329 ERROR  24731 --- [exec-6] c.e.stub.PaymentController  : Failed to parse request body: unexpected EOF"
%%   keyword "ERROR" colored #EF4444, rest of line #EF4444.
%%
%%   NORMAL lines:
%%   "2024-11-01 10:35:45.330  INFO  24731 --- [exec-6] c.e.stub.PaymentController  : Returning 400 Bad Request"
%%   "2024-11-01 10:36:00.000  INFO  24731 --- [Scheduler] c.e.stub.HealthCheck        : Health check OK — uptime 00:00:55"
%%   "2024-11-01 10:36:15.117  INFO  24731 --- [exec-7] c.e.stub.PaymentController  : GET /payment/status → 200 OK"
%%   "2024-11-01 10:36:28.441  INFO  24731 --- [exec-8] c.e.stub.PaymentController  : POST /payment/refund → 200 OK"
%%
%%   Last line ends with a blinking cursor: thin 2×14 px vertical bar, #4F8EF7.
%%
%%   Color coding recap for keywords inside lines:
%%   "INFO"  → #22C55E   "WARN"  → #EAB308   "ERROR" → #EF4444
%%   Timestamps → #4A5568.  Logger names → #64748B.
%%
%% Overall feel: a real Linux terminal embedded inside a polished dark panel.
%% ============================================================


\subsection{Итоги проектирования интерфейса}

В рамках данного раздела разработаны высокодетализированные прототипы шести ключевых экранных форм программного комплекса: общий хром приложения с навигацией (рисунок~\ref{fig:mockup-shell}), главная панель управления Dashboard (рисунок~\ref{fig:mockup-dashboard}), модальное окно регистрации заглушки (рисунок~\ref{fig:mockup-register}), страница настроек с таблицей хостов (рисунок~\ref{fig:mockup-settings}), форма добавления нового хоста (рисунок~\ref{fig:mockup-add-host}) и страница журналов процессов (рисунок~\ref{fig:mockup-logs}).

Все экранные формы выдержаны в единой визуальной системе с тёмной цветовой схемой, акцентным синим цветом (\texttt{\#4F8EF7}) и цветовым кодированием статусов объектов через бейджи и индикаторы. Цветовая система обеспечивает мгновенную оценку состояния инфраструктуры без чтения текстовых подписей.

Разработанные прототипы обосновывают соответствие интерфейса современным требованиям к инструментам класса Operations: минимум навигационных переходов для выполнения целевых задач, низкий порог освоения для инженеров тестирования без специальной веб-разработческой компетенции, а также технологическая независимость от внешних шрифтовых и ресурсных CDN.
